% =====================================================================
%  PARTE DE SARA  --  se incluye desde main.tex con % =====================================================================
%  PARTE DE SARA  --  se incluye desde main.tex con % =====================================================================
%  PARTE DE SARA  --  se incluye desde main.tex con % =====================================================================
%  PARTE DE SARA  --  se incluye desde main.tex con \input{Sara/sara}
%  Contenido: Introducción, Marco Teórico y Tabla de datos
%
%  Reglas para editar sin romper la compilación:
%   * NO escribir \documentclass, \begin{document} ni \end{document} aquí.
%   * Citas: agregar la referencia en ../references.bib y citar con \cite{clave}.
%   * Fórmulas químicas: \ce{CH3NH2}, \ce{(CH3)2NH}, \ce{CH3CONH2}.
%   * Unidades: \si{\gram\per\mole}, \si{\degreeCelsius}, \si{\gram\per\centi\metre\cubed}.
% =====================================================================

\section{Introducción}
\fontsize{12}{14}\selectfont
% TODO (Sara): pegar aquí la redacción de la introducción.
% Sugerencia de hilo argumentativo:
%  1. Importancia de los compuestos nitrogenados (aminas y amidas) en química
%     orgánica, biología y en el ambiente.
%  2. Problema: cómo la estructura (grado de sustitución del N, grupo carbonilo,
%     capacidad de formar puentes de hidrógeno) determina las propiedades físicas.
%  3. Objetivo del trabajo: relacionar estadísticamente propiedades obtenidas en
%     CAS-SciFinder (PM, pf, P.eb, densidad, pKa, momento dipolar, solubilidad).
\textit{[Redacción de la introducción pendiente.]}

Los datos de partida se obtuvieron de CAS-SciFinder \cite{1_Su2022}.
% (Reemplazar esta cita de ejemplo por la referencia real de CAS-SciFinder.)

\section{Marco Teórico}
\fontsize{12}{14}\selectfont
% TODO (Sara): pegar aquí la redacción del marco teórico.
% Temas sugeridos:
%  * Aminas primarias, secundarias y terciarias: \ce{CH3NH2}, \ce{(CH3)2NH}, \ce{(CH3)3N}.
%  * Amidas: acetamida, N-metilacetamida, N,N-dimetilacetamida, urea.
%  * Fuerzas intermoleculares: puentes de hidrógeno N--H$\cdots$N y N--H$\cdots$O,
%    dipolo--dipolo y dispersión de London.
%  * Basicidad (pKa del ácido conjugado) y efecto de la resonancia en amidas.

\subsection{Fuerzas intermoleculares}
\textit{[Redacción pendiente.]}

\subsection{Basicidad y resonancia}
\textit{[Redacción pendiente.]}

\subsection{Tabla de datos de propiedades químicas}
\label{sec:tabla}
% TODO (Sara): completar la tabla con los valores definitivos del Excel
% (Datos/datos_power_bi.xlsx). Una fila por compuesto; usar "--" si falta el dato.
La Tabla~\ref{tab:propiedades} resume las propiedades fisicoquímicas de los
compuestos estudiados, consolidadas desde CAS-SciFinder.

\begin{table}[htbp]
	\centering
	\captionsetup{width=\columnwidth, justification=justified}
	\caption{Propiedades fisicoquímicas de los compuestos estudiados.}
	\label{tab:propiedades}
	\resizebox{\columnwidth}{!}{%
		\begin{tabular}{lccccc}
			\toprule
			\textbf{Compuesto} & \textbf{PM} & \textbf{pf} & \textbf{P.eb} & \textbf{$\rho$} & \textbf{pKa} \\
			 & (\si{\gram\per\mole}) & (\si{\degreeCelsius}) & (\si{\degreeCelsius}) & (\si{\gram\per\centi\metre\cubed}) & \\
			\midrule
			Metilamina            & -- & -- & -- & -- & -- \\
			Dimetilamina          & -- & -- & -- & -- & -- \\
			Trimetilamina         & -- & -- & -- & -- & -- \\
			N-metiletanamina      & -- & -- & -- & -- & -- \\
			Bencidina             & -- & -- & -- & -- & -- \\
			Acetamida             & -- & -- & -- & -- & -- \\
			N-metilacetamida      & -- & -- & -- & -- & -- \\
			N,N-dimetilacetamida  & -- & -- & -- & -- & -- \\
			Urea                  & -- & -- & -- & -- & -- \\
			\bottomrule
		\end{tabular}%
	}
\end{table}

%  Contenido: Introducción, Marco Teórico y Tabla de datos
%
%  Reglas para editar sin romper la compilación:
%   * NO escribir \documentclass, \begin{document} ni \end{document} aquí.
%   * Citas: agregar la referencia en ../references.bib y citar con \cite{clave}.
%   * Fórmulas químicas: \ce{CH3NH2}, \ce{(CH3)2NH}, \ce{CH3CONH2}.
%   * Unidades: \si{\gram\per\mole}, \si{\degreeCelsius}, \si{\gram\per\centi\metre\cubed}.
% =====================================================================

\section{Introducción}
\fontsize{12}{14}\selectfont
% TODO (Sara): pegar aquí la redacción de la introducción.
% Sugerencia de hilo argumentativo:
%  1. Importancia de los compuestos nitrogenados (aminas y amidas) en química
%     orgánica, biología y en el ambiente.
%  2. Problema: cómo la estructura (grado de sustitución del N, grupo carbonilo,
%     capacidad de formar puentes de hidrógeno) determina las propiedades físicas.
%  3. Objetivo del trabajo: relacionar estadísticamente propiedades obtenidas en
%     CAS-SciFinder (PM, pf, P.eb, densidad, pKa, momento dipolar, solubilidad).
\textit{[Redacción de la introducción pendiente.]}

Los datos de partida se obtuvieron de CAS-SciFinder \cite{1_Su2022}.
% (Reemplazar esta cita de ejemplo por la referencia real de CAS-SciFinder.)

\section{Marco Teórico}
\fontsize{12}{14}\selectfont
% TODO (Sara): pegar aquí la redacción del marco teórico.
% Temas sugeridos:
%  * Aminas primarias, secundarias y terciarias: \ce{CH3NH2}, \ce{(CH3)2NH}, \ce{(CH3)3N}.
%  * Amidas: acetamida, N-metilacetamida, N,N-dimetilacetamida, urea.
%  * Fuerzas intermoleculares: puentes de hidrógeno N--H$\cdots$N y N--H$\cdots$O,
%    dipolo--dipolo y dispersión de London.
%  * Basicidad (pKa del ácido conjugado) y efecto de la resonancia en amidas.

\subsection{Fuerzas intermoleculares}
\textit{[Redacción pendiente.]}

\subsection{Basicidad y resonancia}
\textit{[Redacción pendiente.]}

\subsection{Tabla de datos de propiedades químicas}
\label{sec:tabla}
% TODO (Sara): completar la tabla con los valores definitivos del Excel
% (Datos/datos_power_bi.xlsx). Una fila por compuesto; usar "--" si falta el dato.
La Tabla~\ref{tab:propiedades} resume las propiedades fisicoquímicas de los
compuestos estudiados, consolidadas desde CAS-SciFinder.

\begin{table}[htbp]
	\centering
	\captionsetup{width=\columnwidth, justification=justified}
	\caption{Propiedades fisicoquímicas de los compuestos estudiados.}
	\label{tab:propiedades}
	\resizebox{\columnwidth}{!}{%
		\begin{tabular}{lccccc}
			\toprule
			\textbf{Compuesto} & \textbf{PM} & \textbf{pf} & \textbf{P.eb} & \textbf{$\rho$} & \textbf{pKa} \\
			 & (\si{\gram\per\mole}) & (\si{\degreeCelsius}) & (\si{\degreeCelsius}) & (\si{\gram\per\centi\metre\cubed}) & \\
			\midrule
			Metilamina            & -- & -- & -- & -- & -- \\
			Dimetilamina          & -- & -- & -- & -- & -- \\
			Trimetilamina         & -- & -- & -- & -- & -- \\
			N-metiletanamina      & -- & -- & -- & -- & -- \\
			Bencidina             & -- & -- & -- & -- & -- \\
			Acetamida             & -- & -- & -- & -- & -- \\
			N-metilacetamida      & -- & -- & -- & -- & -- \\
			N,N-dimetilacetamida  & -- & -- & -- & -- & -- \\
			Urea                  & -- & -- & -- & -- & -- \\
			\bottomrule
		\end{tabular}%
	}
\end{table}

%  Contenido: Introducción, Marco Teórico y Tabla de datos
%
%  Reglas para editar sin romper la compilación:
%   * NO escribir \documentclass, \begin{document} ni \end{document} aquí.
%   * Citas: agregar la referencia en ../references.bib y citar con \cite{clave}.
%   * Fórmulas químicas: \ce{CH3NH2}, \ce{(CH3)2NH}, \ce{CH3CONH2}.
%   * Unidades: \si{\gram\per\mole}, \si{\degreeCelsius}, \si{\gram\per\centi\metre\cubed}.
% =====================================================================

\section{Introducción}
\fontsize{12}{14}\selectfont
% TODO (Sara): pegar aquí la redacción de la introducción.
% Sugerencia de hilo argumentativo:
%  1. Importancia de los compuestos nitrogenados (aminas y amidas) en química
%     orgánica, biología y en el ambiente.
%  2. Problema: cómo la estructura (grado de sustitución del N, grupo carbonilo,
%     capacidad de formar puentes de hidrógeno) determina las propiedades físicas.
%  3. Objetivo del trabajo: relacionar estadísticamente propiedades obtenidas en
%     CAS-SciFinder (PM, pf, P.eb, densidad, pKa, momento dipolar, solubilidad).
\textit{[Redacción de la introducción pendiente.]}

Los datos de partida se obtuvieron de CAS-SciFinder \cite{1_Su2022}.
% (Reemplazar esta cita de ejemplo por la referencia real de CAS-SciFinder.)

\section{Marco Teórico}
\fontsize{12}{14}\selectfont
% TODO (Sara): pegar aquí la redacción del marco teórico.
% Temas sugeridos:
%  * Aminas primarias, secundarias y terciarias: \ce{CH3NH2}, \ce{(CH3)2NH}, \ce{(CH3)3N}.
%  * Amidas: acetamida, N-metilacetamida, N,N-dimetilacetamida, urea.
%  * Fuerzas intermoleculares: puentes de hidrógeno N--H$\cdots$N y N--H$\cdots$O,
%    dipolo--dipolo y dispersión de London.
%  * Basicidad (pKa del ácido conjugado) y efecto de la resonancia en amidas.

\subsection{Fuerzas intermoleculares}
\textit{[Redacción pendiente.]}

\subsection{Basicidad y resonancia}
\textit{[Redacción pendiente.]}

\subsection{Tabla de datos de propiedades químicas}
\label{sec:tabla}
% TODO (Sara): completar la tabla con los valores definitivos del Excel
% (Datos/datos_power_bi.xlsx). Una fila por compuesto; usar "--" si falta el dato.
La Tabla~\ref{tab:propiedades} resume las propiedades fisicoquímicas de los
compuestos estudiados, consolidadas desde CAS-SciFinder.

\begin{table}[htbp]
	\centering
	\captionsetup{width=\columnwidth, justification=justified}
	\caption{Propiedades fisicoquímicas de los compuestos estudiados.}
	\label{tab:propiedades}
	\resizebox{\columnwidth}{!}{%
		\begin{tabular}{lccccc}
			\toprule
			\textbf{Compuesto} & \textbf{PM} & \textbf{pf} & \textbf{P.eb} & \textbf{$\rho$} & \textbf{pKa} \\
			 & (\si{\gram\per\mole}) & (\si{\degreeCelsius}) & (\si{\degreeCelsius}) & (\si{\gram\per\centi\metre\cubed}) & \\
			\midrule
			Metilamina            & -- & -- & -- & -- & -- \\
			Dimetilamina          & -- & -- & -- & -- & -- \\
			Trimetilamina         & -- & -- & -- & -- & -- \\
			N-metiletanamina      & -- & -- & -- & -- & -- \\
			Bencidina             & -- & -- & -- & -- & -- \\
			Acetamida             & -- & -- & -- & -- & -- \\
			N-metilacetamida      & -- & -- & -- & -- & -- \\
			N,N-dimetilacetamida  & -- & -- & -- & -- & -- \\
			Urea                  & -- & -- & -- & -- & -- \\
			\bottomrule
		\end{tabular}%
	}
\end{table}

%  Contenido: Introducción, Marco Teórico y Tabla de datos
%
%  Reglas para editar sin romper la compilación:
%   * NO escribir \documentclass, \begin{document} ni \end{document} aquí.
%   * Citas: agregar la referencia en ../references.bib y citar con \cite{clave}.
%   * Fórmulas químicas: \ce{CH3NH2}, \ce{(CH3)2NH}, \ce{CH3CONH2}.
%   * Unidades: \si{\gram\per\mole}, \si{\degreeCelsius}, \si{\gram\per\centi\metre\cubed}.
% =====================================================================

\section{Introducción}
\fontsize{12}{14}\selectfont
% TODO (Sara): pegar aquí la redacción de la introducción.
% Sugerencia de hilo argumentativo:
%  1. Importancia de los compuestos nitrogenados (aminas y amidas) en química
%     orgánica, biología y en el ambiente.
%  2. Problema: cómo la estructura (grado de sustitución del N, grupo carbonilo,
%     capacidad de formar puentes de hidrógeno) determina las propiedades físicas.
%  3. Objetivo del trabajo: relacionar estadísticamente propiedades obtenidas en
%     CAS-SciFinder (PM, pf, P.eb, densidad, pKa, momento dipolar, solubilidad).
\textit{[Redacción de la introducción pendiente.]}

Los datos de partida se obtuvieron de CAS-SciFinder \cite{1_Su2022}.
% (Reemplazar esta cita de ejemplo por la referencia real de CAS-SciFinder.)

\section{Marco Teórico}
\fontsize{12}{14}\selectfont
% TODO (Sara): pegar aquí la redacción del marco teórico.
% Temas sugeridos:
%  * Aminas primarias, secundarias y terciarias: \ce{CH3NH2}, \ce{(CH3)2NH}, \ce{(CH3)3N}.
%  * Amidas: acetamida, N-metilacetamida, N,N-dimetilacetamida, urea.
%  * Fuerzas intermoleculares: puentes de hidrógeno N--H$\cdots$N y N--H$\cdots$O,
%    dipolo--dipolo y dispersión de London.
%  * Basicidad (pKa del ácido conjugado) y efecto de la resonancia en amidas.

\subsection{Fuerzas intermoleculares}
\textit{[Redacción pendiente.]}

\subsection{Basicidad y resonancia}
\textit{[Redacción pendiente.]}

\subsection{Tabla de datos de propiedades químicas}
\label{sec:tabla}
% TODO (Sara): completar la tabla con los valores definitivos del Excel
% (Datos/datos_power_bi.xlsx). Una fila por compuesto; usar "--" si falta el dato.
La Tabla~\ref{tab:propiedades} resume las propiedades fisicoquímicas de los
compuestos estudiados, consolidadas desde CAS-SciFinder.

\begin{table}[htbp]
	\centering
	\captionsetup{width=\columnwidth, justification=justified}
	\caption{Propiedades fisicoquímicas de los compuestos estudiados.}
	\label{tab:propiedades}
	\resizebox{\columnwidth}{!}{%
		\begin{tabular}{lccccc}
			\toprule
			\textbf{Compuesto} & \textbf{PM} & \textbf{pf} & \textbf{P.eb} & \textbf{$\rho$} & \textbf{pKa} \\
			 & (\si{\gram\per\mole}) & (\si{\degreeCelsius}) & (\si{\degreeCelsius}) & (\si{\gram\per\centi\metre\cubed}) & \\
			\midrule
			Metilamina            & -- & -- & -- & -- & -- \\
			Dimetilamina          & -- & -- & -- & -- & -- \\
			Trimetilamina         & -- & -- & -- & -- & -- \\
			N-metiletanamina      & -- & -- & -- & -- & -- \\
			Bencidina             & -- & -- & -- & -- & -- \\
			Acetamida             & -- & -- & -- & -- & -- \\
			N-metilacetamida      & -- & -- & -- & -- & -- \\
			N,N-dimetilacetamida  & -- & -- & -- & -- & -- \\
			Urea                  & -- & -- & -- & -- & -- \\
			\bottomrule
		\end{tabular}%
	}
\end{table}
